% W4i40_Experiments.tex
% DFD 17-Agent Research Programme — Iteration 40 (i40)
% Agents 9 (Experimental Proposals), 10 (DFD vs Everything),
%         13 (BMV Landscape), 14 (Clock Network)
% Iteration 9/20 of Quantum Gravity Cycle (i32-i51)
% Date: 2026-03-22

\documentclass[11pt,a4paper]{article}
\usepackage[utf8]{inputenc}
\usepackage{amsmath,amssymb,amsfonts,amsthm}
\usepackage{hyperref}
\usepackage{booktabs}
\usepackage{longtable}
\usepackage{geometry}
\usepackage{xcolor}
\geometry{margin=2.5cm}

\newtheorem{theorem}{Theorem}
\newtheorem{proposition}{Proposition}
\newtheorem{lemma}{Lemma}
\newtheorem{corollary}{Corollary}
\newtheorem{definition}{Definition}
\newtheorem{remark}{Remark}

\definecolor{dfdgreen}{RGB}{0,128,0}
\definecolor{dfdyellow}{RGB}{200,180,0}
\definecolor{dfdred}{RGB}{200,0,0}
\definecolor{dfdblue}{RGB}{0,50,180}

\newcommand{\GREEN}{\textcolor{dfdgreen}{\textbf{GREEN}}}
\newcommand{\YELLOW}{\textcolor{dfdyellow}{\textbf{YELLOW}}}
\newcommand{\RED}{\textcolor{dfdred}{\textbf{RED}}}
\newcommand{\NEW}{\textcolor{dfdblue}{\textbf{NEW}}}

\title{DFD i40: Experimental Proposals and Theory Comparison\\[6pt]
\large Agents 9 (Experiments), 10 (DFD vs All), 13 (BMV Labs), 14 (Clocks)\\[6pt]
\normalsize Iteration 9/20 of Quantum Gravity Cycle (i32--i51)}
\author{DFD 17-Agent Research Programme}
\date{2026-03-22}

\begin{document}
\maketitle
\tableofcontents
\newpage

%=============================================================================
\section{Agent 9: Experimental Proposals --- Ranked}
\label{sec:agent9}
%=============================================================================

\subsection{Mandate}

Rank ALL DFD-testable experiments by: discrimination power $\times$ feasibility / cost. Identify the top 5 with realistic budgets and timelines.

\subsection{Key References}

\begin{itemize}
\item Bose et al.\ (2017), PRL 119, 240401 --- BMV proposal
\item Carney, Stamp \& Taylor (2019), CQG 36, 034001 --- tabletop QG experiments
\item MAQRO-PF White Paper (2025), arXiv:2512.01777
\item Menlo Systems (2025), BACON network achievement
\item LIGO-Virgo-KAGRA (2024), O4 observing run
\item Adelberger et al.\ (2009), ARNPS 59, 141 --- Eot-Wash torsion balance
\item McGaugh et al.\ (2016), PRL 117, 201101 --- radial acceleration relation
\item Desmond et al.\ (2024), MNRAS --- wide binary tests of gravity
\end{itemize}

\subsection{Complete Experiment Ranking}

\begin{center}
\begin{longtable}{clcccc}
\toprule
\textbf{Rank} & \textbf{Experiment} & \textbf{Discrim.} & \textbf{Feasib.} & \textbf{Cost} & \textbf{Score} \\
 & & (1--10) & (1--10) & (relative) & $D\times F/C$ \\
\midrule
\endhead
1 & BMV (space-based) & 10 & 4 & 5 & 8.0 \\
2 & GW echoes (O5/O6) & 8 & 7 & 3 & 18.7 \\
3 & Wide binary kinematics & 7 & 9 & 1 & 63.0 \\
4 & Optical clock network & 6 & 8 & 2 & 24.0 \\
5 & Galaxy RAR precision & 8 & 8 & 2 & 32.0 \\
6 & Torsion balance (Eot-Wash) & 7 & 9 & 1 & 63.0 \\
7 & Atom interferometer (MAGIS) & 7 & 6 & 4 & 10.5 \\
8 & CMB (SymBoltz prediction) & 9 & 5 & 0.5 & 90.0 \\
9 & Cluster dynamics (eROSITA) & 6 & 7 & 2 & 21.0 \\
10 & MAQRO-PF (space) & 8 & 3 & 8 & 3.0 \\
\bottomrule
\end{longtable}
\end{center}

\textbf{Note:} Score = Discrimination $\times$ Feasibility / Cost (higher = better). Cost is relative: 1 = use existing data, 10 = \$1B+ mission.

\subsection{Top 5 Detailed Proposals}

\subsubsection{Proposal 1: CMB Power Spectrum Prediction (Score: 90)}

\textbf{Description:} Run DFD through SymBoltz.jl and predict the full CMB $C_\ell$ spectrum. Compare to Planck 2020 data.

\textbf{Discrimination:} 9/10 --- the CMB is the single most constraining dataset for any cosmological theory. If DFD matches the first three peaks without CDM, it is revolutionary.

\textbf{Feasibility:} 5/10 --- requires implementing the nonlinear MOND constraint in a linear Boltzmann solver (the subject of Agents 1--2). The QUMOND linearization may not capture all the physics.

\textbf{Cost:} 0.5 --- purely computational. No new experiments needed. Existing Planck data.

\textbf{Timeline:} 3--6 months after SymBoltz implementation is complete.

\textbf{Budget:} \$0 (computational resources only; graduate student time).

\textbf{What would confirm DFD:} $\chi^2_{\rm DFD}/\chi^2_{\Lambda\rm CDM} < 2$ for $\ell < 2000$.

\textbf{What would falsify DFD:} $R_{31}^{\rm DFD} < 0.30$ or $> 0.55$ (outside $3\sigma$ of Planck).

\subsubsection{Proposal 2: Wide Binary Kinematics (Score: 63)}

\textbf{Description:} Use Gaia DR4 data to measure orbital velocities of wide binary stars ($a > 5000$ AU) where internal accelerations $< a_0$. DFD predicts MOND-enhanced velocities.

\textbf{Discrimination:} 7/10 --- directly tests the MOND interpolation function at accessible accelerations. Ongoing controversy (Chae 2023 vs.\ Banik et al.\ 2024).

\textbf{Feasibility:} 9/10 --- Gaia DR4 expected 2026--2027 with improved proper motions.

\textbf{Cost:} 1 --- uses existing/upcoming public data.

\textbf{Timeline:} 1--2 years after Gaia DR4.

\textbf{Budget:} \$0 (data analysis only).

\textbf{DFD prediction:} At separations $> 7$ kAU, relative velocities exceed Newtonian by factor $\sqrt{\nu}$. The transition should follow $\mu(g/a_0) = g/g_N$ with $a_0 = 1.2\times10^{-10}$ m/s$^2$.

\textbf{Key issue:} The external field effect. Wide binaries are embedded in the Galactic gravitational field ($g_{\rm ext} \sim 1.8\times10^{-10}$ m/s$^2 \sim 1.5\,a_0$). In MOND, the EFE reduces the enhancement. DFD must specify its EFE prediction precisely.

\subsubsection{Proposal 3: Torsion Balance / Fifth Force (Score: 63)}

\textbf{Description:} Eot-Wash-style torsion balance experiments testing short-range deviations from Newtonian gravity. DFD predicts $G_{\rm eff} = G_N\nu(g_N/a_0)$, which varies at different acceleration scales.

\textbf{Discrimination:} 7/10 --- directly tests $G_{\rm eff}$ vs.\ $G_N$.

\textbf{Feasibility:} 9/10 --- Eot-Wash group operates continuously.

\textbf{Cost:} 1 --- incremental improvement to existing apparatus.

\textbf{DFD prediction:} At laboratory accelerations ($g \sim 10^{-7}$ m/s$^2$ for mg masses at cm separations), $y = g/a_0 \sim 10^3$, giving $\nu \approx 1.001$. This is a $0.1\%$ deviation from Newtonian gravity --- at the edge of current precision.

\textbf{Challenge:} The EFE from Earth's gravity ($g = 9.8$ m/s$^2 \gg a_0$) should suppress MOND effects. Laboratory tests may not access the MOND regime.

\subsubsection{Proposal 4: Galaxy RAR Precision (Score: 32)}

\textbf{Description:} Use JWST, VLT/MUSE, and SKA precursor surveys to measure the radial acceleration relation (RAR) with $< 5\%$ precision across $10^4$ galaxies spanning $10^{-13}$ to $10^{-8}$ m/s$^2$.

\textbf{Discrimination:} 8/10 --- the RAR is DFD's strongest current prediction. Existing data already shows $< 0.13$ dex scatter.

\textbf{Feasibility:} 8/10 --- surveys are ongoing (WALLABY, MHONGOOSE, etc.).

\textbf{Cost:} 2 --- piggybacks on existing large surveys.

\textbf{Timeline:} 2--4 years for sufficient statistics.

\textbf{DFD prediction:} The RAR follows $g_{\rm obs} = \nu(g_{\rm bar}/a_0)\,g_{\rm bar}$ with \emph{zero} scatter from a universal $\mu$ function. Any intrinsic scatter $> 0.05$ dex would indicate additional physics (dark matter halos or variations in $a_0$).

\subsubsection{Proposal 5: Optical Clock Network (Score: 24)}

\textbf{Description:} Use the BACON network and European fiber-linked optical clocks (NIST, JILA, PTB, SYRTE, NPL) to test gravitational redshift with $10^{-18}$ fractional frequency precision.

\textbf{Discrimination:} 6/10 --- DFD modifies the gravitational potential, which affects clock rates. The modification is $\delta\nu/\nu = \nu(g/a_0)\Phi/c^2$ vs.\ $\Phi/c^2$ in GR.

\textbf{Feasibility:} 8/10 --- BACON achieved $8\times10^{-18}$ uncertainty in 2025. Six-country network demonstrated.

\textbf{Cost:} 2 --- uses existing infrastructure.

\textbf{Timeline:} 1--3 years for dedicated measurement campaigns.

\textbf{DFD prediction:} For clocks at different altitudes on Earth ($g \approx 9.8$ m/s$^2 \gg a_0$), $\nu \approx 1$ and DFD matches GR to $< 10^{-19}$. The DFD signature appears only at accelerations $\lesssim a_0$, requiring free-falling or space-based clocks.

\textbf{Specific opportunity:} Clocks on the International Space Station experience $g \sim 8.7$ m/s$^2$, still $\gg a_0$. Clocks on a highly eccentric orbit might sample different acceleration regimes, but the variation is always $\gg a_0$.

\textbf{HONEST ASSESSMENT:} Clock networks cannot test DFD's unique predictions unless placed in a microgravity environment where $g \lesssim a_0$. This requires a dedicated space mission.

\subsection{Summary: Experiment Priority Matrix}

\begin{center}
\begin{tabular}{lcccl}
\toprule
\textbf{Experiment} & \textbf{Tests} & \textbf{Timeline} & \textbf{Budget} & \textbf{Status} \\
\midrule
SymBoltz CMB & Cosmology & 3--6 mo & \$0 & \textbf{DO NOW} \\
Wide binaries & MOND $\mu(x)$ & 1--2 yr & \$0 & Wait for Gaia DR4 \\
Torsion balance & $G_{\rm eff}$ & 1 yr & \$100K & Incremental \\
Galaxy RAR & MOND universality & 2--4 yr & \$0 & Ongoing surveys \\
Clock network & Gravitational redshift & 1--3 yr & \$500K & BACON ready \\
BMV (space) & Quantum gravity & 5--10 yr & \$500M+ & Concept phase \\
GW echoes & $D_0$ & 3--5 yr & \$0 & Wait for O5/O6 \\
MAQRO-PF & QG decoherence & 5--8 yr & \$200M & Phase A \\
\bottomrule
\end{tabular}
\end{center}

\textbf{Grade: B+}

%=============================================================================
\section{Agent 10: DFD vs Everything --- Head-to-Head}
\label{sec:agent10}
%=============================================================================

\subsection{Mandate}

Compare DFD against ALL competitors. Where does DFD win? Where does it lose? Be brutal.

\subsection{Key References}

\begin{itemize}
\item Skordis \& Z{\l}o\'{s}nik (2021), PRL 127, 161302 --- AeST
\item Berezhiani \& Khoury (2015), PRD 92, 103510 --- superfluid DM
\item Hu \& Sawicki (2007), PRD 76, 064004 --- $f(R)$ gravity
\item Bekenstein (2004), PRD 70, 083509 --- TeVeS
\item Milgrom (2023), PRD 108, 063505 --- MOND review
\item Blanchet et al.\ (2024), JCAP 2024, 040 --- relativistic khronon
\item Durakovi\'{c} et al.\ (2025), arXiv:2503.11174 --- mimetic gravity
\item Planck Collaboration (2020), A\&A 641, A6
\end{itemize}

\subsection{Comparison Matrix}

\begin{center}
\begin{longtable}{lccccccc}
\toprule
\textbf{Observable} & \textbf{DFD} & \textbf{$\Lambda$CDM} & \textbf{AeST} & \textbf{SFDM} & \textbf{$f(R)$} & \textbf{TeVeS} & \textbf{RMOND} \\
\midrule
\endhead
\multicolumn{8}{l}{\textbf{Galaxy Scale}} \\
\midrule
Rotation curves & \GREEN & \YELLOW & \GREEN & \GREEN & \RED & \GREEN & \GREEN \\
Tully--Fisher & \GREEN & \RED & \GREEN & \GREEN & \RED & \GREEN & \GREEN \\
RAR & \GREEN & \YELLOW & \GREEN & \GREEN & \RED & \GREEN & \GREEN \\
EFE & \GREEN & N/A & \GREEN & \YELLOW & N/A & \GREEN & \GREEN \\
\midrule
\multicolumn{8}{l}{\textbf{Cluster Scale}} \\
\midrule
Cluster masses & \YELLOW & \GREEN & \YELLOW & \YELLOW & \GREEN & \RED & \YELLOW \\
Bullet Cluster & \RED & \GREEN & \YELLOW & \YELLOW & \GREEN & \RED & \YELLOW \\
Lensing & \YELLOW & \GREEN & \GREEN & \YELLOW & \GREEN & \YELLOW & \YELLOW \\
\midrule
\multicolumn{8}{l}{\textbf{Cosmological}} \\
\midrule
CMB $C_\ell$ & \YELLOW & \GREEN & \GREEN & \YELLOW & \GREEN & \RED & \YELLOW \\
BAO & \YELLOW & \GREEN & \GREEN & \YELLOW & \GREEN & \RED & \YELLOW \\
$\sigma_8$ tension & \GREEN & \RED & \YELLOW & \YELLOW & \YELLOW & --- & \YELLOW \\
$H_0$ tension & \YELLOW & \RED & \YELLOW & \YELLOW & \YELLOW & --- & \YELLOW \\
LSS $P(k)$ & \YELLOW & \GREEN & \GREEN & \YELLOW & \GREEN & \RED & \YELLOW \\
\midrule
\multicolumn{8}{l}{\textbf{Quantum/Fundamental}} \\
\midrule
Born rule & \GREEN & N/A & N/A & N/A & N/A & N/A & N/A \\
Decoherence & \GREEN & N/A & N/A & N/A & N/A & N/A & N/A \\
BMV prediction & \GREEN & \YELLOW & N/A & N/A & N/A & N/A & N/A \\
GW echoes & \GREEN & N/A & \YELLOW & N/A & \YELLOW & N/A & N/A \\
Dark energy & \YELLOW & \GREEN & \YELLOW & \YELLOW & \GREEN & \YELLOW & \YELLOW \\
\midrule
\multicolumn{8}{l}{\textbf{Theoretical}} \\
\midrule
Parameters & 1 & 6 & 4+ & 3 & 2 & 5+ & 3 \\
Simplicity & \GREEN & \GREEN & \YELLOW & \YELLOW & \GREEN & \RED & \YELLOW \\
Stability & \GREEN & \GREEN & \GREEN & \YELLOW & \GREEN & \RED & \YELLOW \\
\bottomrule
\end{longtable}
\end{center}

Legend: \GREEN = good agreement, \YELLOW = partial/uncertain, \RED = tension/failure. SFDM = Superfluid Dark Matter. RMOND = Relativistic MOND (generic). N/A = not applicable (theory doesn't address this).

\subsection{Detailed Head-to-Head Comparisons}

\subsubsection{DFD vs $\Lambda$CDM}

\textbf{DFD wins:}
\begin{itemize}
\item Galaxy rotation curves: DFD predicts them from baryonic mass alone. $\Lambda$CDM requires fitting dark matter halo parameters for each galaxy.
\item Tully--Fisher relation: DFD predicts $v^4 = G_N M a_0$ exactly. $\Lambda$CDM must invoke baryonic feedback to explain the tight relation.
\item RAR: DFD gives $g_{\rm obs} = \nu(g_{\rm bar}/a_0)g_{\rm bar}$ with zero free parameters. $\Lambda$CDM requires fine-tuned baryonic feedback.
\item $\sigma_8$ tension: DFD naturally suppresses small-scale power (MOND enhances large-scale more than small-scale in some regimes).
\item Born rule / quantum gravity: $\Lambda$CDM says nothing about quantum foundations.
\item Parameter count: DFD has 1 parameter ($a_0$) vs.\ $\Lambda$CDM's 6.
\end{itemize}

\textbf{DFD loses:}
\begin{itemize}
\item CMB: $\Lambda$CDM fits all $\ell$ from 2 to 2500. DFD is untested (YELLOW).
\item Bullet Cluster: The spatial offset between mass (lensing) and baryons (X-ray) is natural in $\Lambda$CDM but requires explanation in DFD.
\item BAO: $\Lambda$CDM matches BAO measurements precisely. DFD's modified sound horizon is untested.
\item Large-scale structure: $\Lambda$CDM has well-tested N-body simulations. DFD has none.
\item Inflation: Both need external inflation, but $\Lambda$CDM's inflationary predictions (nearly scale-invariant spectrum) are well-tested.
\end{itemize}

\textbf{Verdict: DFD is better at galaxies, worse at cosmology (pending SymBoltz). Quantum gravity gives DFD unique predictions that $\Lambda$CDM cannot match.}

\subsubsection{DFD vs AeST}

\textbf{AeST strengths:}
\begin{itemize}
\item Fits full CMB $C_\ell$ spectrum (demonstrated by Skordis \& Zlosnik 2021).
\item Gravitational lensing matches observations.
\item Stable, well-studied theory with Hamiltonian formulation (6 DOF confirmed).
\item BAO compatible.
\end{itemize}

\textbf{DFD strengths:}
\begin{itemize}
\item Fewer parameters (1 vs.\ 4+).
\item Quantum gravity predictions (Born rule, BMV, decoherence).
\item Simpler action (single scalar, no vector field needed).
\item GW echo predictions.
\end{itemize}

\textbf{AeST weaknesses (DFD opportunities):}
\begin{itemize}
\item AeST requires a vector field with 4 parameters ($K_B$, $K_S$, etc.).
\item AeST's MOND function must be tuned separately from the cosmological function.
\item AeST says nothing about quantum mechanics.
\end{itemize}

\textbf{DFD weaknesses (vs AeST):}
\begin{itemize}
\item DFD has not demonstrated CMB fit. AeST has.
\item DFD may need a vector field anyway (if CMB fails without one).
\item DFD's quantum gravity claims are unproven.
\end{itemize}

\textbf{Verdict: AeST is the benchmark. DFD must match its CMB fit or explain why it differs. DFD's advantage is quantum gravity + simplicity.}

\subsubsection{DFD vs Superfluid Dark Matter (SFDM)}

\textbf{SFDM concept:} Dark matter forms a superfluid on galaxy scales, with phonon-mediated force mimicking MOND. On cluster/cosmological scales, it behaves as standard CDM.

\textbf{SFDM strengths:}
\begin{itemize}
\item Natural transition between MOND (galaxy) and CDM (cosmology).
\item Explains Bullet Cluster (DM halos exist but are superfluid).
\item CMB compatible (acts as CDM at recombination).
\end{itemize}

\textbf{DFD strengths:}
\begin{itemize}
\item No dark matter particle needed (Occam's razor).
\item Quantum gravity predictions.
\item Fewer parameters.
\item SFDM requires a new particle + a phase transition + fine-tuned phonon coupling.
\end{itemize}

\textbf{Verdict: SFDM is clever but adds complexity (new particle + phase transition). DFD is simpler but must explain clusters and CMB.}

\subsubsection{DFD vs $f(R)$ Gravity}

\textbf{$f(R)$ strengths:}
\begin{itemize}
\item Explains cosmic acceleration without $\Lambda$.
\item Well-studied, implemented in multiple Boltzmann solvers (MGCAMB, hi\_class).
\item Solar system constraints satisfied via chameleon mechanism.
\end{itemize}

\textbf{DFD strengths:}
\begin{itemize}
\item $f(R)$ does NOT produce MOND phenomenology. It cannot explain galaxy rotation curves or the RAR without dark matter.
\item DFD addresses dark matter AND dark energy AND quantum gravity. $f(R)$ only addresses dark energy.
\item DFD has fewer parameters.
\end{itemize}

\textbf{Verdict: Not really competitors. $f(R)$ is a dark energy theory; DFD is a dark matter + quantum gravity theory. They address different problems.}

\subsubsection{DFD vs TeVeS}

\textbf{TeVeS is effectively dead:}
\begin{itemize}
\item Gravitational wave speed: TeVeS predicted $c_{\rm GW} \neq c$. GW170817/GRB170817A measured $|c_{\rm GW}/c - 1| < 10^{-15}$. TeVeS ruled out.
\item CMB: TeVeS failed to reproduce the third peak (the same challenge DFD faces).
\item Replaced by AeST, which satisfies $c_{\rm GW} = c$.
\end{itemize}

\textbf{DFD vs TeVeS: DFD wins on GW speed, similar challenges on CMB.}

\subsection{The Brutal Honest Assessment}

\begin{center}
\begin{tabular}{lcc}
\toprule
\textbf{Theory} & \textbf{DFD wins} & \textbf{DFD loses} \\
\midrule
$\Lambda$CDM & Galaxies, QG, parameters & CMB, clusters, LSS \\
AeST & Simplicity, QG & CMB (demonstrated fit) \\
SFDM & Simplicity, no new particle & Clusters, CMB transition \\
$f(R)$ & Galaxies, scope & (different problems) \\
TeVeS & GW speed, stability & (TeVeS is dead) \\
\bottomrule
\end{tabular}
\end{center}

\textbf{DFD's unique selling point:} No other theory attempts to unify MOND, quantum gravity, and the Born rule. If the CMB works out, DFD has no competitor.

\textbf{DFD's biggest risk:} CMB. If SymBoltz.jl shows $R_{31}$ is wrong, DFD needs a vector field (becoming AeST-like) or fails.

\textbf{Grade: B+}

%=============================================================================
\section{Agent 13: BMV Experimental Landscape}
\label{sec:agent13}
%=============================================================================

\subsection{Mandate}

Map ALL BMV-type experiments worldwide. Who is doing what? When? What sensitivity? Could they detect DFD's 2200$\times$ enhancement?

\subsection{Key References}

\begin{itemize}
\item Bose et al.\ (2017), PRL 119, 240401
\item Marletto \& Vedral (2017), PRL 119, 240402
\item Carney, Stamp \& Taylor (2019), CQG 36, 034001
\item Matsumura et al.\ (2024), nanodiamond BMV analysis
\item Christodoulou et al.\ (2025), arXiv:2506.21122 --- BMV without spacetime superposition
\item Krisnanda et al.\ (2020), npj Quantum Information 6, 12
\item Marshman, Mazumdar \& Bose (2020), PRD 101, 016014
\item van de Kamp et al.\ (2020), PRA 101, 062301
\end{itemize}

\subsection{Global BMV/GIE Experimental Groups}

\begin{center}
\begin{longtable}{llccl}
\toprule
\textbf{Group} & \textbf{Institution} & \textbf{Mass} & \textbf{Status} & \textbf{Timeline} \\
 & & (kg) & & \\
\midrule
\endhead
Bose group & UCL & $10^{-14}$ & Theory/design & 2028--2030 \\
Vedral group & Oxford & $10^{-14}$ & Theory/design & 2028--2030 \\
Aspelmeyer group & Vienna & $10^{-15}$ & Lab prototype & 2027--2029 \\
Arndt group & Vienna & $10^{-18}$ & Interferometry & 2026--2028 \\
Romero-Isart group & Innsbruck & $10^{-17}$ & Levitation & 2027--2029 \\
Geraci group & Northwestern & $10^{-13}$ & Sensor dev. & 2027--2030 \\
Moore group & Yale & $10^{-12}$ & Force sensing & 2026--2028 \\
Giessibl group & Regensburg & $10^{-15}$ & AFM-based & Concept \\
MAQRO consortium & ESA & $10^{-15}$ & Phase A & 2030+ \\
Chinese groups & USTC, PKU & Various & Early stage & 2028+ \\
\bottomrule
\end{longtable}
\end{center}

\subsection{Sensitivity Analysis: Can They Detect DFD?}

\subsubsection{Standard GR BMV Signal}

For masses $m = 10^{-14}$ kg at separation $d = 200\,\mu$m:
\begin{equation}
\phi_{\rm GR} = \frac{G_N m^2 T}{\hbar d} = \frac{6.67\times10^{-11}\times10^{-28}\times T}{10^{-34}\times2\times10^{-4}} = 3.3\times10^{-1}\,T\;\text{[rad]}
\end{equation}

For $\phi_{\rm GR} = \pi/2$: $T \approx 4.8$ s. This requires coherence time $> 5$ s for a $10^{-14}$ kg mass. Current decoherence times for levitated nanoparticles: $\sim 10^{-3}$ s. Gap: $\sim 5000\times$.

\subsubsection{DFD-Enhanced Signal (Ground-Based)}

On Earth: $g_{\rm ext} = 9.8$ m/s$^2 \gg a_0$, so $\nu \approx 1$. \textbf{No enhancement.}

\subsubsection{DFD-Enhanced Signal (Free-Fall)}

In free-fall or microgravity ($g_{\rm ext} < 10^{-6}$ m/s$^2$):
\begin{equation}
y = \frac{G_N m}{a_0 d^2} = \frac{6.67\times10^{-11}\times10^{-14}}{1.2\times10^{-10}\times4\times10^{-8}} = 1.4\times10^{-7}
\end{equation}
\begin{equation}
\nu \approx \sqrt{1/y} = 2680
\end{equation}

Required coherence time: $T_{\rm DFD} = T_{\rm GR}/2680 = 1.8$ ms. \textbf{This IS within current technology!}

But: free-fall on Earth lasts $< 10$ s (drop tower), and the residual acceleration is $\sim 10^{-6}$ m/s$^2$, which gives $y_{\rm ext} = g_{\rm ext}/a_0 \sim 10^4$, so $\nu \approx 1$. \textbf{Still no enhancement.}

The EFE question is critical:
\begin{itemize}
\item If DFD's EFE is like MOND's: only the self-gravity of the test masses matters for the \emph{relative} potential, and the external field affects only the \emph{overall} acceleration. Then the enhancement applies even on Earth (because the self-gravity potential difference $\Delta V$ between superposition branches is in the MOND regime regardless of background).
\item If DFD's EFE is absolute: the total acceleration field determines $\nu$, and on Earth $\nu \approx 1$.
\end{itemize}

\textbf{THIS IS THE SINGLE MOST IMPORTANT THEORETICAL QUESTION FOR BMV.}

\subsubsection{The EFE Resolution}

In the QUMOND formulation:
\begin{equation}
\nabla^2\Phi_{\rm true} = \nabla\cdot\left[\nu\!\left(\frac{|\nabla\Phi_N|}{a_0}\right)\nabla\Phi_N\right]
\end{equation}

The total Newtonian potential $\Phi_N$ includes \emph{all} sources (test masses + Earth + Sun + ...). The gradient $\nabla\Phi_N$ is dominated by the external field. So at a given point:
\begin{equation}
|\nabla\Phi_N| \approx g_{\rm ext} + \delta g_{\rm self}
\end{equation}

where $\delta g_{\rm self} \ll g_{\rm ext}$ is the self-gravity of the test masses. Then:
\begin{equation}
\nu\!\left(\frac{g_{\rm ext} + \delta g}{a_0}\right) \approx \nu\!\left(\frac{g_{\rm ext}}{a_0}\right) + \nu'\!\left(\frac{g_{\rm ext}}{a_0}\right)\frac{\delta g}{a_0}
\end{equation}

For $g_{\rm ext} \gg a_0$: $\nu(g_{\rm ext}/a_0) \approx 1$ and $\nu' \approx 0$. The self-gravity enhancement is suppressed.

\textbf{CONCLUSION:} In QUMOND (and by extension DFD), the BMV enhancement is suppressed by the external field on Earth. The experiment requires microgravity with $g_{\rm ext} \lesssim a_0 = 1.2\times10^{-10}$ m/s$^2$.

\textbf{Required platform:} LISA Pathfinder-class drag-free satellite ($g \sim 10^{-15}$ m/s$^2$).

\subsection{Agent 13 Summary}

\textbf{Grade: B}

The BMV landscape is active but no group is close to the GR-level signal ($\sim 5000\times$ gap in coherence time). DFD's enhancement could close this gap \emph{but only in microgravity}. Ground-based experiments cannot access the MOND regime for self-gravitating test masses due to the EFE.

\textbf{Recommendation:} Propose a space-based BMV experiment on a LISA Pathfinder follow-on mission. Alternatively, investigate whether \emph{any} formulation of DFD allows partial EFE shielding.

%=============================================================================
\section{Agent 14: Optical Clock Network Tests of DFD}
\label{sec:agent14}
%=============================================================================

\subsection{Mandate}

Can a global optical clock network test DFD? What precision is needed? Assess BACON and European fiber links.

\subsection{Key References}

\begin{itemize}
\item Menlo Systems (2025), BACON network links --- sub-$8\times10^{-18}$ uncertainty
\item Optica (2025), unprecedented six-country optical clock network
\item Lisdat et al.\ (2016), Nature Photonics 10, 662 --- optical clock networks
\item Grotti et al.\ (2018), Nature Physics 14, 437 --- geodesy with clocks
\item Takamoto et al.\ (2020), Nature Photonics 14, 411 --- cm-level geodesy
\item Adelberger et al.\ (2009), ARNPS 59, 141 --- equivalence principle tests
\item Will (2014), LRR 17, 4 --- testing general relativity
\end{itemize}

\subsection{What BACON Achieved (2025)}

The Boulder Atomic Clock Optical Network achieved:
\begin{itemize}
\item Frequency comparison uncertainty: $< 8\times10^{-18}$ for Al$^+$/Sr, Al$^+$/Yb, Yb/Sr ratios
\item Links: 3.6 km fiber (NIST--JILA) + 1.5 km free-space optical link
\item Six-country international comparison: 10 clocks, 38 simultaneous frequency ratios
\item Precision: $100\times$ better than satellite techniques
\end{itemize}

\subsection{DFD Predictions for Clock Networks}

\subsubsection{Gravitational Redshift}

The fractional frequency shift between two clocks at gravitational potentials $\Phi_1$ and $\Phi_2$:

\textbf{GR:}
\begin{equation}
\frac{\Delta\nu}{\nu} = \frac{\Phi_1 - \Phi_2}{c^2}
\end{equation}

\textbf{DFD:}
\begin{equation}
\frac{\Delta\nu}{\nu} = \frac{\Phi_1^{\rm DFD} - \Phi_2^{\rm DFD}}{c^2}
\end{equation}

where $\Phi^{\rm DFD} = \nu(g_N/a_0)\Phi_N$. On Earth's surface, $g \approx 9.8$ m/s$^2$, $y = g/a_0 = 8.2\times10^{10}$, so:
\begin{equation}
\nu(8.2\times10^{10}) = 1 + 2.4\times10^{-11} + \mathcal{O}(y^{-2})
\end{equation}

The DFD correction to the gravitational redshift on Earth:
\begin{equation}
\frac{\delta(\Delta\nu/\nu)}{\Delta\nu/\nu} \approx \frac{\nu(y) - 1}{1} \approx 2.4\times10^{-11}
\end{equation}

For a height difference $\Delta h = 1$ m: $\Delta\nu/\nu = g\Delta h/c^2 = 1.09\times10^{-16}$.

DFD correction: $\delta(\Delta\nu/\nu) = 2.4\times10^{-11}\times1.09\times10^{-16} = 2.6\times10^{-27}$.

\textbf{This is 9 orders of magnitude below current clock precision.} Clock networks cannot test DFD on Earth.

\subsubsection{Can Anything Be Done?}

\textbf{Option 1: Solar System Tests.}
Clock on a spacecraft falling toward the Sun. At $r = 1$ AU, $g_\odot = 6\times10^{-3}$ m/s$^2$, $y = 5\times10^7$. DFD correction: $\sim 4\times10^{-8}$ fractional. Still too small relative to the redshift itself.

\textbf{Option 2: Gravitational potential difference between Earth and deep space.}
A clock in deep space (far from any mass) would experience the cosmic acceleration field. If $g_{\rm cosmic} \lesssim a_0$, the MOND correction would be significant. But measuring the redshift between Earth and deep space requires a spacecraft mission.

\textbf{Option 3: Indirect test via $D_0$.}
The DFD coupling $D_0 = 0.017$ affects the fine-structure constant through gravitational loop corrections. The predicted shift: $\delta\alpha/\alpha \sim D_0^2 G_N m_e^2/(4\pi\hbar c) \sim 10^{-60}$. Completely undetectable.

\textbf{Option 4: Clock comparison in different gravitational environments.}
Compare clocks at different latitudes (different $g$ due to Earth's oblateness + rotation). The variation $\delta g/g \sim 10^{-3}$ means $\delta\nu/\nu \sim 10^{-3}\times2.4\times10^{-11} \sim 10^{-14}$... times the redshift, giving $\sim 10^{-30}$. Undetectable.

\subsection{Negative Result: Clocks Cannot Test DFD (on Earth)}

\textbf{Grade: C}

\textbf{HONEST CONCLUSION:} Optical clock networks, despite their extraordinary precision ($8\times10^{-18}$), cannot test DFD's unique predictions because:
\begin{enumerate}
\item Earth's gravitational acceleration is $10^{10}\times$ above $a_0$
\item The MOND correction to $G_{\rm eff}$ is $\sim 10^{-11}$
\item This translates to a $\sim 10^{-27}$ correction to the redshift, far below $10^{-18}$ precision
\end{enumerate}

\textbf{What WOULD work:}
\begin{itemize}
\item A clock on a drag-free satellite in deep space ($g < a_0$): DFD predicts a measurable shift in the gravitational redshift for masses in the MOND regime.
\item Clock comparisons between different positions in the Milky Way halo (different $g$ relative to $a_0$): requires interstellar probes.
\end{itemize}

\textbf{Clock networks are excellent for testing GR but cannot discriminate DFD from GR.} The DFD-specific signatures require accessing the MOND acceleration regime, which is not possible on Earth or in low Earth orbit.

\subsection{What Clocks CAN Do for DFD}

While clocks cannot directly test the MOND interpolation function, they can:

\begin{enumerate}
\item \textbf{Constrain DFD parameter space:} If DFD predicts any gravitational anomaly at $g \gg a_0$ (e.g., from $D_0$ coupling), clocks could detect it. But currently, DFD predicts zero anomaly in this regime.

\item \textbf{Test the universality of free fall:} DFD's constraint field couples to total mass-energy, not just rest mass. This could produce composition-dependent effects at $\sim D_0^2$ level. Current WEP tests: $\eta < 10^{-15}$ (MICROSCOPE). DFD prediction: $\eta \sim D_0^2 \sim 3\times10^{-4}$... wait, that seems too large.

Actually, let us compute carefully. The DFD effective gravitational coupling for a body of mass $m$:
\begin{equation}
G_{\rm eff}(m) = G_N\left(1 + D_0\frac{\delta\rho_{\rm binding}}{m/R^3}\right)
\end{equation}

where $\delta\rho_{\rm binding}$ is the gravitational binding energy contribution. For laboratory masses, $\delta\rho_{\rm binding}/\rho \sim G_N m/(Rc^2) \sim 10^{-25}$. So:
\begin{equation}
\eta \sim D_0 \times 10^{-25} \sim 10^{-27}
\end{equation}

This is far below MICROSCOPE's sensitivity. No WEP violation detectable.
\end{enumerate}

%=============================================================================
\section{File Summary}
\label{sec:summary}
%=============================================================================

\begin{center}
\begin{tabular}{lcl}
\toprule
\textbf{Agent} & \textbf{Grade} & \textbf{Key Finding} \\
\midrule
9 (Experiments) & B+ & CMB computation is \#1 priority; wide binaries \#2 \\
10 (DFD vs All) & B+ & DFD unique in QG; CMB is make-or-break \\
13 (BMV Labs) & B & EFE kills ground-based BMV; need space \\
14 (Clocks) & C & Cannot test DFD on Earth ($g \gg a_0$) \\
\bottomrule
\end{tabular}
\end{center}

\textbf{Critical findings:}
\begin{enumerate}
\item The external field effect (EFE) is the single biggest obstacle for laboratory tests of DFD's quantum predictions. All MOND-regime experiments require microgravity ($g < a_0$).
\item The CMB prediction (SymBoltz.jl) is the cheapest, fastest, and most discriminating test. It should be done IMMEDIATELY.
\item Wide binary kinematics (Gaia DR4) is the best near-term observational test.
\item DFD's unique advantage over all competitors is its quantum gravity sector. No other modified gravity theory predicts the Born rule or BMV enhancement.
\item The BMV experiment needs a space platform --- no ground-based experiment can access the MOND regime for self-gravitating test masses.
\end{enumerate}

\end{document}
